\clearpage
\seccionInforme{5. Apreciación de la herramienta}
\phantomsection\addcontentsline{toc}{section}{5. Apreciación de la herramienta}

La valoración se basa en lo que el equipo observó al construir y probar MatriLab con React Native y Expo, no en afirmaciones generales sobre el framework.

\vspace{0.2cm}
\phantomsection\addcontentsline{toc}{subsection}{5.1 Características favorables}
\colabtitle{5.1 Características favorables}

\subsubtitulo{A. Interfaces nativas reales a partir de un solo código}

Cada componente del núcleo se convierte en una vista nativa del sistema [2]. En la práctica, el \texttt{Switch} de la configuración es el interruptor de Android, el teclado y la hoja de compartir son los del teléfono, el \texttt{Alert} es el diálogo del sistema y \texttt{native-stack} usa las transiciones y el gesto de retroceso nativos. Con ese mismo código, sin una sola línea específica, la aplicación se exportó a la web (Figura 21). Las diferencias por plataforma que sí hubo que escribir fueron mínimas y explícitas: el \texttt{behavior} de \texttt{Keyboard\allowbreak{}Avoiding\allowbreak{}View} y el \texttt{android\_\allowbreak{}ripple} de \texttt{Pressable}. Para un proyecto universitario que debe llegar a estudiantes con Android y con iPhone, esto reduce a la mitad el trabajo frente al desarrollo nativo separado.

\subsubtitulo{B. Ciclo de desarrollo rápido y ecosistema JavaScript}

Con Expo Go no fue necesario compilar proyectos de Android Studio: el emulador cargaba el paquete desde Metro y \textit{Fast Refresh} aplicaba cada cambio de estilo en uno o dos segundos conservando el estado de la pantalla, lo que permitió ajustar márgenes, columnas y colores viendo el resultado al instante. El ecosistema de JavaScript aportó piezas maduras: TypeScript detectó errores de navegación entre pantallas antes de ejecutarlas, \texttt{react-native-qrcode-svg} resolvió el QR en una línea y Jest probó las reglas de negocio en un segundo. El mismo lenguaje sirvió para la aplicación, las pruebas y el script de capturas.

\vspace{0.2cm}
\phantomsection\addcontentsline{toc}{subsection}{5.2 Características desfavorables}
\colabtitle{5.2 Características desfavorables}

\subsubtitulo{A. Dependencia de la compatibilidad entre versiones}

Una aplicación React Native depende de una cadena de versiones que deben coincidir: Expo SDK, React Native, React, cada librería con código nativo (\texttt{react-native-screens}, \texttt{react-native-svg}, \texttt{expo-camera}) y la versión de Expo Go instalada en el teléfono, que solo abre proyectos de su mismo SDK. Fue obligatorio instalar las dependencias con \texttt{npx expo install} para no romper esa cadena. El ritmo de cambios también se notó: el \texttt{Safe\allowbreak{}Area\allowbreak{}View} del núcleo figura como obsoleto [16], la Nueva Arquitectura cambió el comportamiento interno de varias librerías y TypeScript 6 dejó de incluir tipos globales por defecto, lo que obligó a corregir \texttt{tsconfig.\allowbreak{}json} (cambio 6). En un proyecto que dure varios semestres, actualizar dependencias requiere tiempo dedicado.

\subsubtitulo{B. Layout limitado a Flexbox y diferencias entre plataformas}

React Native solo ofrece Flexbox y posición absoluta: no hay CSS Grid ni \textit{media queries}. La cuadrícula del horario tuvo que calcularse a mano con \texttt{use\allowbreak{}Window\allowbreak{}Dimensions} y \texttt{position: 'absolute'}, y la rejilla adaptable de laboratorios obligó a volver a montar la \texttt{Flat\allowbreak{}List} con \texttt{key} porque \texttt{num\allowbreak{}Columns} no puede cambiar en caliente. Algunos componentes se comportan distinto en cada sistema (\texttt{Keyboard\allowbreak{}Avoiding\allowbreak{}View}, sombras, ripple) y hay que probar en ambos. Además, la depuración de lo que ocurre del lado nativo exige conocer las herramientas de Android: los permisos de cámara de Expo Go, el teclado del emulador y la rotación tuvieron que resolverse con adb.

\vspace{0.2cm}
\phantomsection\addcontentsline{toc}{subsection}{5.3 Valoración general}
\colabtitle{5.3 Valoración general}

\begin{table}[H]
\centering\small
\caption{Resumen de la apreciación de React Native con Expo}
\renewcommand{\arraystretch}{1.25}
\begin{tabular}{|>{\raggedright\arraybackslash}p{\dimexpr 0.280\linewidth-2\tabcolsep-2\arrayrulewidth\relax}|>{\raggedright\arraybackslash}p{\dimexpr 0.160\linewidth-2\tabcolsep-2\arrayrulewidth\relax}|>{\raggedright\arraybackslash}p{\dimexpr 0.560\linewidth-2\tabcolsep-2\arrayrulewidth\relax}|}
\hline
\rowcolor{colabpyblue} \color{white}\textbf{Aspecto} & \color{white}\textbf{Valoración} & \color{white}\textbf{Evidencia en el laboratorio} \\ \hline
Componentes propios de la plataforma & Favorable & Switch, Alert, Share, teclado y pila de navegación nativos en Android. \\ \hline
\rowcolor{colabpybg} Productividad & Favorable & Fast Refresh, Expo Go, TypeScript y pruebas del dominio en 1 s. \\ \hline
Mantenimiento de versiones & Desfavorable & Cadena SDK–RN–librerías–Expo Go; APIs obsoletas; ajuste de tsconfig. \\ \hline
\rowcolor{colabpybg} Sistema de layout & Desfavorable & Sin Grid ni media queries; horario y rejilla calculados a mano. \\ \hline
\end{tabular}
\end{table}

